\documentclass[11pt]{article}
\usepackage{amsmath,amssymb,amsthm,hyperref}
\newtheorem{lemma}{Lemma}
\newtheorem{proposition}{Proposition}
\title{Balanced geometric clusters: exact spectrum and Hermite cost}
\author{Research checkpoint 007}
\date{31 October 2026}
\begin{document}
\maketitle

\paragraph{Scope.}
We analyze a particular replacement of consecutive packet coordinates:
unions of complete low-bit clusters at geometric offsets. We charge the
standard full Hermite readout with triangular multiplicities and its
uniform per-node precision, followed by unrestricted feasible-mode
enumeration. The resulting candidate bounds exclude an exponent gain
from these replacements within that method. They are not complexity
lower bounds for all interpolation functionals, support-aware algorithms,
or ECDLP. No full C38 contraction or improved complete ECDLP bound is
claimed. The ultimate square-root-in-$b$ goal remains unchanged.

\section{Exact derivative structure}
For $n\ge4$, let
\[
 X_n=\{0,1,2,4,\ldots,2^{n-2}\},\quad D_n=X_n-X_n,
 \quad E_n=D_n\cup(D_n-1).
\]
Write $m=n-2$ and define
\[
 d=|D_n|=m^2+m+3,\quad
 c=(n-3)(n-4)+1,\quad
 e=|E_n|=2(n^2-5n+9)=d+c.
\]
Every odd point of $D_n$ has both adjacent even points in $D_n$.
Its even part is $2D_{n-1}$. Hence
\begin{equation}
 E_n=2D_{n-1}\ \cup\ (2D_{n-1}-1).
 \label{eq:coarse}
\end{equation}
In particular, the $c$ points added to $D_n$ are all odd.

Put $B=2^u$ and use packet coordinates
\[
 X=B X_n+\{0,\ldots,B-1\},\qquad N=|X|=nB.
\]
Every residue modulo $2^v$, $v\le u$, occurs $nB/2^v$ times.
The distinct difference grid is
\begin{equation}
 \Gamma=BD_n\ \cup\ \bigcup_{r=1}^{B-1}(BE_n+r),
 \qquad |\Gamma|=Be-c.
 \label{eq:grid}
\end{equation}
All quantities here refer to distinct roots, with no multiplicity
silently imported from representations of a difference.

For a finite integer set $H$, let
$V_H(t)=\sum_{s\in H\setminus\{t\}}v_2(t-s)$.
Writing $V_0=V_{D_n}(0)=((n-2)^3+5(n-2))/3$, one has
\[
 V_{E_n}(0)=V_0,\qquad
 \min_{t\in E_n}V_{E_n}(t)=n^2-2n-3.
\]
The first identity follows because the added roots are odd. For the
second, \eqref{eq:coarse} gives
$V_{E_n}(2s)=V_{E_n}(2s-1)=|D_{n-1}|-1+V_{D_{n-1}}(s)$;
checkpoint 006 proved $\min V_{D_{n-1}}=3n-11$.

\begin{lemma}
For $\tau=Bt+r\in\Gamma$, $1\le r<B$, and $i=v_2(r)$,
\begin{equation}
 V_\Gamma(\tau)-V_\Gamma(0)
 =c(u-i)+V_{E_n}(t)-V_0.
 \label{eq:derivative}
\end{equation}
For $r=0$ the difference is $V_{D_n}(t)-V_0$.
\end{lemma}
\begin{proof}
At zero, the same-residue contribution is $u(d-1)+V_0$.
Other residues contribute $e(B-1-u)$, since
$\sum_{r=1}^{B-1}v_2(r)=B-1-u$.
At a nonzero residue $r$, the same-residue contribution is
$u(e-1)+V_{E_n}(t)$. If every residue had $e$ roots, the other-residue
contribution would again be $e(B-1-u)$: a full residue system modulo
$2^u$ has constant derivative valuation. Residue zero has $c$ fewer
roots, subtracting $ci$. Subtract the zero-node expression.
\end{proof}

Let $v\le u$ low-bit levels be retained, so all nodes with
$v_2(\tau)<v$ are used. Set
\[
 A_n=V_0-(n^2-2n-3).
\]
Every retained constant coefficient is integral exactly when
\begin{equation}
 c(u-v+1)\ge A_n.
 \label{eq:admissible}
\end{equation}
This is necessary because $W_{xy}(\tau)$ is nonzero when $y-x=\tau$;
its valuation is \eqref{eq:derivative}. Every grid node is an actual
coordinate difference. It is sufficient for constant terms by the same
formula. The local scale exponent $\ell=u+n-1$ is sufficient for all
coefficients: it bounds every pair-difference valuation in $\Gamma$.

For $m=n-2$,
\[
 \frac{A_n}{c}=\frac m3-\frac{4m-9}{3c}
 \ge\frac{m-1}{3}.
\]
The inequality uses $(m-3)(m-4)\ge0$ for integer $m\ge2$.
Thus \eqref{eq:admissible} implies
\begin{equation}
 n\le3(u-v+1)+3\le3u+3.
 \label{eq:cluster-size}
\end{equation}
For example, $n=8,u=2$ allows $v=1$, but not $v=2$.
An actual retained node at the latter level has constant valuation $-16$,
even though the central node has positive constant valuation.

\section{The normalized mode spectrum}
Consider the central odd node $\tau=1$ and the even coordinate block
\[
 2\bigl(b X_n+\{0,\ldots,b-1\}\bigr),\qquad b=B/2.
\]
Enumerate it by increasing cluster and then increasing low coordinate.
For rank $j=ab+t$, $0\le a<n$, $0\le t<b$, the normalized Cauchy
diagonal increment is
\begin{equation}
 \Delta_j=4j-2s_2(t)+a(a-3).
 \label{eq:mode}
\end{equation}
Indeed the valuation sum over each earlier full cluster is
$b-1+v_2(x_a-x_{a'})$, and the partial current cluster contributes
$v_2(t!)=t-s_2(t)$. The geometric cluster prefix contributes
$a(a-1)/2$. Doubling the total valuation sum after accounting for the
outer factor two gives \eqref{eq:mode}.

This order also gives integral normalized triangular factors. The
geometric cluster order minimizes the prefix valuation sum among
remaining clusters. Within a cluster, divisibility of a product of
$t$ consecutive integers by $t!$ proves the corresponding minimum
property. Thus every remaining row's numerator valuation is at least
the pivot's; all odd-node Cauchy denominators are units. Each normalized
triangular factor has unit diagonal. Direct elimination and an
independent greedy-minimum certificate check these facts in the code.

The common crossing weight is
\[
 \delta=\ell+cu,
\]
by \eqref{eq:derivative} at $t=0$. Hence
\begin{equation}
 w_j=\ell+cu+\Delta_j
 \le\ell+cu+4j+n^2.
 \label{eq:weight-upper}
\end{equation}
Compared with consecutive coordinates, the extra diagonal valuation is
$a(a-3)+2s_2(a)$. Under \eqref{eq:cluster-size}, this is only polynomial
in $u$, while each cluster contains $2^{u-1}$ points. The readout
calculation below makes the resulting enumeration limitation precise.

\section{The full Hermite readout cost}
Use multiplicities $q_\tau=h_0(v-v_2(\tau))$ at retained nodes and
let the local substitution be $T=\tau+2^\ell Z$.
For any retained-node set $\mathcal N$, define
\[
 E_\tau(0)=\prod_{s\in\mathcal N\setminus\{\tau\}}
       \left(\frac{-s}{\tau-s}\right)^{q_s}.
\]
The coefficient of the highest local jet in the standard full Hermite
evaluation functional at zero is exactly
\begin{equation}
 E_\tau(0)\left(\frac{-\tau}{2^\ell}\right)^{q_\tau-1}.
 \label{eq:highest}
\end{equation}
There is no sum in this particular coefficient that could cancel its
denominator. Uniform per-node precision for this functional must charge
its negative valuation, unless additional coefficient relations or
structure are used. Such additional methods are outside this result.
The code independently verifies the complete functional on every
monomial through its interpolation degree in small cases.

For $h=u-i$, $z=v-i$, put
\[
 \mathcal A(h,z)=3\cdot2^h-2^{h-z}(z+3)-z(h+1).
\]
On the positive central residue set $1,\ldots,B-1$, restricted to the
retained levels, the negative valuation of $E_\tau(0)$ is
$h_0\mathcal A(h,z)$.
To see this, sum separately over nodes with valuation less than, equal
to, and greater than $i$. The lower valuations cancel. The remaining
sum is
\[
 h_0\left[z\sum_{j=1}^{h-1}j2^{h-j-1}
       -\sum_{j=1}^{z-1}j(z-j)2^{h-j-1}\right],
\]
which equals the displayed expression.

\subsection{All coarse nodes retained at each active level}
For $\tau=Bt+r$, the full-grid negative valuation is
\begin{equation}
 -v_2(E_\tau(0))
 =h_0\left[e\mathcal A(h,z)
       +z\bigl((e-1)h+V_{E_n}(t)\bigr)\right].
 \label{eq:full-loss}
\end{equation}
Different residues contribute $e$ copies of the central-residue sum;
the other $e-1$ roots of the same residue contribute the second term.
At $\tau=1$ call this loss $L_0$. The highest-jet loss is
\[
 L_0+(h_0v-1)\ell.
\]
The number of retained jets is
\[
 h_0 eB(v-1+2^{-v}).
\]
To interpolate a deformation of degree $M(|\Gamma|-1)$ it must exceed
that degree. Therefore $h_0\ge M/(2v)$, using $c+1\le e$ and $B\ge2$.
For $u\ge32$, \eqref{eq:full-loss} gives
\begin{equation}
 L_0\ge\tfrac12 h_0eB\ge\frac{Mn^2B}{8v}.
 \label{eq:precision-lower}
\end{equation}
Here $(v+3)/2^v\le2$, $v\le u$, $2^u$ dominates the discarded
quadratic terms, and $e\ge n^2/2$.

Assume for clarity $M/(nB)$ full crossing packets. Choose the first
\[
 k=\left\lfloor\frac{nB}{16}
                 \min\{1,\sqrt{n/v}\}\right\rfloor
\]
modes from one orientation in each packet. For $u\ge32$ and
\eqref{eq:cluster-size}, $cu+n^2\le k$.
One elementary uniform certificate is
$2^u\ge144u(u+1)^3+16u$, true at $u=32$ and preserved by induction.
Thus the total non-$\ell$ part of all selected weights is at most
\[
 3\frac{M}{nB}k^2
 \le\frac{3}{256}\frac{Mn^2B}{v}<L_0.
\]
Their number $K=Mk/(nB)$ is at most $M/16$, whereas
$q_1=h_0v\ge M/2$. The additional $K\ell$ cost is covered by the
$(q_1-1)\ell$ part of \eqref{eq:highest}, for any nonnegative local
scale exponent. Every binary vector on these $K$ modes is therefore
feasible in both degree and precision. Moreover
\begin{equation}
 K\ge\frac M{32}\min\{1,\sqrt{n/v}\}.
 \label{eq:candidates}
\end{equation}
With a fixed positive fraction of crossing packets, as in balanced
group packing, only this absolute constant changes.

Consequently unrestricted enumeration contains
$2^{\Omega(M\min\{1,\sqrt{n/v}\})}$ candidates. Since
$v\le u\le\log_2 M$, this does not improve the old
$M/\sqrt{\log M}$ enumeration scale; growing $n$ makes this bound
larger. This is a candidate-universe result, not a claim that all
candidates contribute to a particular tensor network.

\subsection{Only positive central nodes retained}
Now use only $\tau\in\{1,\ldots,B-1\}$ with $v_2(\tau)<v$.
Every such node has nonnegative constant valuation, for arbitrary $n$.
The coarse factor $e$ disappears from the jet count but the deformation
degree is unchanged. Hence
\[
 h_0\ge\frac{Me}{2v},\qquad q_1\ge Me/2,
 \qquad L_0=h_0\mathcal A(u,v).
\]
For $u\ge2$, $\mathcal A(u,v)\ge B/4$: it increases with $v$, with
minimum $B-u-1$, and $u=2$ is the endpoint case.
Therefore $L_0\ge Mn^2B/(16v)$.

If $n<v$ and $u\ge32$, the same prefix argument as above applies,
since $n\le u$ and $3/256<1/16$.
If $n\ge v$, choose any fixed $M/128$ modes, taking $M$ divisible
by $128nB$ to avoid rounding. Their non-$\ell$ weights are at most
\[
 \frac M{128}\bigl(n^2(u+1)+2nB\bigr)<L_0,
\]
using $v(u+1)\le4B$ and $n\ge v$. Their degree and $\ell$ costs
are covered by $q_1$. Thus this regime already has $2^{M/128}$
feasible candidates.

For $u=v=1$ the only retained node is $1$, so the full Taylor readout
has $q_1-1=M(|\Gamma|-1)$ and highest loss $(q_1-1)\ell$.
The scale $\ell=0$ is nonintegral: the leading coefficient of each
$W$ polynomial is $1/C_\Gamma'(0)$. At the required jet length this
coefficient is present. Hence $\ell\ge1$. Now
$|\Gamma|-1\ge n^2$ and the maximum mode weight is at most
$\ell+3n^2$, again admitting $M/128$ freely chosen binary modes.
The finitely bounded cases $u<32,n<v$ have bounded packet parameters
and do not create a new asymptotic saving as $M\to\infty$.

\paragraph{What remains open.}
Both conclusions use the standard full Hermite functional and the
unrestricted catalogue of feasible modes. When the number of jets
exceeds the deformation degree, another interpolation functional can
use the extra linear relations; coefficient-dependent precision or
support-aware contraction may also behave differently. No lower bound
against those alternatives is asserted. Arbitrary selective coarse
node sets and nontriangular profiles are not covered either.

\paragraph{Validation.}
Thirty finite fixtures verify exact derivative formulas, integral
normalized factors and their mode weights. Independent elimination,
full monomial interpolation identities, a negative constant-term
control, and large-parameter rational/integer certificates are retained
by the reproducibility launcher. The large certificates evaluate closed
formulas and inequalities; they do not construct enormous packets or
claim experimental speedups on those inputs.
\end{document}
